\section{Turbulent CFD validation}\label{sec:turbulent-validation}

\subsection{Grid sensitivity and streamwise development}

The turbulent grid study in Figure~\ref{fig:turbulent-grid} gives smaller pairwise differences than the laminar study. Medium-to-fine changes are $0.24\%$, $0.13\%$, and $0.37\%$ for radial, axial, and combined refinement, respectively. These values support the use of the retained grids for model comparison, while remaining conditional on the chosen observables.

\begin{figure}[htbp]
    \centering
    \includegraphics[width=\linewidth]{turbulent_grid_independence.pdf}
    \caption{Turbulent-flow grid study: radial velocity, axial pressure, and combined refinement.}
    \label{fig:turbulent-grid}
\end{figure}

The correlation $L_e=1.359DRe^{1/4}$ gives $2.249$~m. The full-profile $0.1\%$ criterion is first satisfied at $z=9.85$~m, close to the final interior station. Figure~\ref{fig:turbulent-development} explains the difference: the centreline velocity reaches an approximately stable level earlier, whereas the full radial profile continues to change weakly. The result illustrates the sensitivity of a turbulent entry length to its definition.

\begin{figure}[htbp]
    \centering
    \includegraphics[width=0.84\linewidth]{turbulent_flow_development.pdf}
    \caption{Centreline-velocity development and axial pressure decrease for the standard $k$--$\varepsilon$ case.}
    \label{fig:turbulent-development}
\end{figure}

\subsection{Resolved mean and turbulence fields}

Figure~\ref{fig:turbulent-structure} collects the fully developed mean velocity, eddy-viscosity ratio, turbulent kinetic energy, and dissipation rate. The mean profile is flatter than the laminar parabola. Turbulent quantities increase towards the wall region, consistent with shear production, while the centreline symmetry is retained.

\begin{figure}[htbp]
    \centering
    \includegraphics[width=0.94\linewidth]{turbulent_structure_profiles.pdf}
    \caption{Radial structure of the standard $k$--$\varepsilon$ solution in the developed region.}
    \label{fig:turbulent-structure}
\end{figure}

The wall stress gives $u_\tau=0.03654$~m/s. Figure~\ref{fig:wall-scaling} compares the computed inner-scaled velocity with the viscous and logarithmic reference laws. The maximum plotted $y^+$ is approximately $1699$ because it corresponds to the distance from the wall to the pipe axis; it is not the first-cell wall coordinate. The profile crosses the expected logarithmic range but the near-wall assessment remains limited by the available cell-centre sampling.

\begin{figure}[htbp]
    \centering
    \includegraphics[width=0.70\linewidth]{turbulent_wall_scaling.pdf}
    \caption{Wall-scaled axial velocity for the standard $k$--$\varepsilon$ case.}
    \label{fig:wall-scaling}
\end{figure}

\subsection{Model comparison}

Figure~\ref{fig:turbulent-profiles} compares the four numerical profiles with two experimental series. All closures reproduce the main shape. The standard model gives the lowest normalised profile RMSE, $0.0158$, followed by RNG ($0.0166$), realisable ($0.0195$), and low-Re ($0.0214$).

\begin{figure}[htbp]
    \centering
    \includegraphics[width=0.76\linewidth]{turbulent_pipe_validation.pdf}
    \caption{Turbulent velocity profiles and experimental measurements at $Re_b=75{,}000$.}
    \label{fig:turbulent-profiles}
\end{figure}

\begin{table}[htbp]
\centering
\caption{Turbulence-model validation metrics.}
\label{tab:turbulent-models}
\begin{tabular}{lrrr}
\toprule
Model & $f$ & $|f-f_{\mathrm{exp}}|/f_{\mathrm{exp}}$ [\%] & Profile RMSE \\
\midrule
Standard $k$--$\varepsilon$ & 0.018984 & 0.082 & 0.01584 \\
RNG $k$--$\varepsilon$ & 0.018403 & 3.143 & 0.01659 \\
Realisable $k$--$\varepsilon$ & 0.017853 & 6.035 & 0.01948 \\
Low-Re $k$--$\varepsilon$ & 0.021688 & 14.148 & 0.02144 \\
\bottomrule
\end{tabular}
\end{table}

The experimental Darcy factor is $0.0190$. The smooth-pipe Prandtl and Haaland values are $0.01912$ and $0.01895$, respectively. Figure~\ref{fig:friction-comparison} shows that the standard model is also closest to all three references. The low-Re configuration overpredicts wall friction in the present implementation; this is consistent with a mismatch between a wall-resolving formulation and the available near-wall resolution rather than with an intrinsic failure of the closure.

\subsection{Low-Reynolds-number resolution sensitivity}

The low-Reynolds-number case was examined separately because its friction-factor discrepancy is larger than that of the other closures. Four radial meshes, with $N_y=35$, 45, 60, and 75 cells, were compared using the developed-region velocity profile at $z/L\simeq0.8$. The successive relative $L_\infty$ changes are $1.10\%$, $0.55\%$, and $3.15\%$. The final change is not smaller than the preceding one; the sequence therefore does not establish monotonic convergence or grid independence for this observable.

\begin{figure}[htbp]
    \centering
    \includegraphics[width=0.82\linewidth]{low_re_grid_sensitivity.pdf}
    \caption{Developed-region low-Reynolds-number profiles for four radial meshes.}
    \label{fig:low-re-grid}
\end{figure}

The three solver exports available on the $N_y=45$ mesh permit a complementary iteration-sensitivity check. From 500 to 750 iterations, the relative $L_\infty$ changes are $0.96\%$ in axial velocity, $1.89\%$ in pressure, $4.58\%$ in turbulent kinetic energy, $6.54\%$ in dissipation rate, and $12.30\%$ in eddy viscosity. From 750 to 1000 iterations, they decrease to $0.23\%$, $0.33\%$, $1.19\%$, $3.03\%$, and $4.05\%$, respectively. The mean-flow fields change less than the turbulence quantities at the final iteration interval, but the latter have not reached a demonstrated asymptote.

\begin{figure}[htbp]
    \centering
    \includegraphics[width=0.84\linewidth]{low_re_iteration_sensitivity.pdf}
    \caption{Relative field changes between successive low-Reynolds-number solver exports.}
    \label{fig:low-re-iterations}
\end{figure}

An additional comparison varies the radial-spacing exponent while keeping the nominal radial cell count fixed at 11. The profile changes are $0.66\%$ between $p=-1.5$ and $p=-1.8$, and $0.37\%$ between $p=-1.8$ and $p=-2.0$. This supports limited sensitivity to the tested spacing choices, but it does not resolve the non-monotonic cell-count study or the remaining turbulence-field iteration changes.

\begin{figure}[htbp]
    \centering
    \includegraphics[width=0.78\linewidth]{power_y_stretching_sensitivity.pdf}
    \caption{Low-Reynolds-number profiles for three radial-spacing exponents at fixed nominal cell count.}
    \label{fig:low-re-spacing}
\end{figure}

\begin{figure}[htbp]
    \centering
    \includegraphics[width=0.80\linewidth]{turbulent_friction_factor_comparison.pdf}
    \caption{Darcy friction factors from the RANS solutions, experiment, and smooth-pipe correlations.}
    \label{fig:friction-comparison}
\end{figure}
